% !TeX root = ../main.tex

% ==================================================
% A.12. 등록형 데코레이터와 FastAPI
% ==================================================

\subsection{등록형 데코레이터와 FastAPI}

Registration Decorator · Decorator Factory · Route Registry · Endpoint · Dependency Injection · Function Introspection

\subsubsection{이전 학습과의 연결}

\zsh{A.11}에서는 클래스 내부에 정의된 함수에 데코레이터를 적용할 때 두 가지 메커니즘을 구별했다.

첫째, 데코레이터 적용은 함수 정의문이 실행되는 과정에서 원본 함수 객체를 다른 객체로 변환할 수 있다.

둘째, 디스크립터 바인딩은 이후 클래스나 인스턴스를 통해 속성을 조회할 때 그 객체의 접근 결과를 결정한다.

이번에는 데코레이터의 또 다른 활용 방식인 등록형 데코레이터(Registration Decorator)를 학습할 것이다.

지금까지 우리가 주로 사용한 데코레이터는 다음과 같은 구조였다.

\begin{dhkcode}
@trace
def calculate(x):
    return x * 2
\end{dhkcode}

\zsh{trace}는 원본 함수 객체를 전달받아 새로운 wrapper 함수 객체를 반환하였다.

그런데 다음 코드를 살펴보자.

\begin{dhkcode}
@app.get("/items/{item_id}")
def get_item(item_id: int):
    return {"item_id": item_id}
\end{dhkcode}

이것은 FastAPI에서 사용하는 경로 처리 데코레이터(Path Operation Decorator)다.

겉으로 보면 지금까지 학습한 데코레이터와 같지만, 목적이 다르다.

\zsh{@app.get(...)}은 원본 함수에 로깅이나 성능 측정 기능을 덧붙이기 위해 wrapper를 생성하는 것이 아니다.

특정 HTTP 요청을 처리할 함수 객체를 FastAPI 애플리케이션에 등록하는 것이 주된 목적이다.

이 차이를 이해하면 데코레이터를 단순한 함수 변환 도구가 아니라, 프로그램의 구성 정보를 선언하고 연결하는 수단으로 바라볼 수 있다.

\subsubsection{이번 학습의 핵심 질문}

이번 학습에서는 다음 질문을 해결할 것이다.

\begin{dhksmry}
1. 등록형 데코레이터는 일반적인 wrapper 데코레이터와 무엇이 다를까?

2. @app.get("/tiems")에서 app.get(...)은 어떤 객체를 반환할까?

3. 함수 객체는 언제 생성되고, 언제 FastAPI에 등록될까?

4. 등록 이후 함수 이름과 FastAPI 내부의 엔드포인트 참조는 어떤 관계일까?

5. HTTP 요청이 들어오면 Python 함수 객체는 어떤 과정을 거쳐 호출될까?

6. FastAPI는 함수의 매개변수와 타입 힌트를 어떻게 분석할까?

7. @app.get과 @trace의 적용 순서를 바꾸면 어떤 차이가 발생할까?

8. 동기 함수와 비동기 함수에 데코레이터를 적용할 때 무엇을 주의해야 할까?
\end{dhksmry}

핵심은 등록 시점(Registration Time)과 요청 처리 시점(Request Time)을 명확히 구별하는 것이다.

\subsubsection{등록형 데코레이터의 정확한 정의}

등록형 데코레이터는 전달받은 객체를 레지스트리(Registry)와 같은 외부 관리 구조에 등록하는 부작용을 수행하는 데코레이터다.

등록 과정이 끝나면 원본 객체를 그대로 반환할 수도 있고, 다른 객체를 반환할 수도 있다.

따라서 등록형 데코레이터의 본질은 wrapper를 생성하지 않는다는 데 있지 않다.

전달받은 객체와 외부 관리 구조 사이에 등록 관계를 형성한다는 데 있다.

다음 두 구현을 비교해 보자.

\begin{dhknote}
\begin{dhksmry}[centered]
Wrapper 데코레이터
\end{dhksmry}
\end{dhknote}

\begin{dhkcode}
def trace(func):
    def wrapper(*args, **kwargs):
        print("TRACE")
        return func(*args, **kwargs)

    return wrapper
\end{dhkcode}

함수 호출을 중개하는 새 함수 객체를 반환한다.

\begin{dhknote}
\begin{dhksmry}[centered]
등록형 데코레이터
\end{dhksmry}
\end{dhknote}

\begin{dhkcode}
registry = {}

def register(func):
    registry[func.__name__] = func
    return func
\end{dhkcode}

함수 객체를 외부 자료구조에 저장하고 원본을 반환한다.

두 번째 예제에서 중요한 문장은 다음이다.

\begin{dhkcode}
registry[func.__name__] = func
\end{dhkcode}

이 문장은 새로운 wrapper를 생성하지 않는다.

대신 \zsh{registry} 딕셔너리에 함수 객체에 대한 참조를 저장한다.

그리고 다음 문장은

\begin{dhkcode}
return func
\end{dhkcode}

원본 함수 객체를 그대로 반환한다.

따라서 데코레이션 이후 함수 이름이 참조하는 객체의 동일성은 유지된다.

하지만 외부 레지스트리에는 새로운 등록 정보가 추가된다.

이것이 등록형 데코레이터의 가장 기본적인 형태다.

\subsubsection{가장 단순한 등록형 데코레이터 구현}

먼저 FastAPI없이 Python만으로 등록형 데코레이터를 만들어 보자.

\begin{dhkcode}
registry = {}


def register(func):
    print("register:", func.__name__)

    registry[func.__name__] = func

    return func


@register
def greet(name):
    print("Hello", name)
    return name.upper()


print("definition complete")

print(registry["greet"] is greet)

result = registry["greet"]("Kim")

print("result:", result)
\end{dhkcode}

예상 출력:

\begin{dhkzsh}
register: greet
definition complete
True
Hello Kim
result: KIM
\end{dhkzsh}

이 결과를 객체 생성과 바인딩 순서에 따라 분석해 보자.

\begin{dhknote}
\begin{dhksmry}[centered]
1. 함수 정의문 실행
\end{dhksmry}
\end{dhknote}

Python이 다음 함수 정의문에 도달한다.

\begin{dhkcode}
@register
def greet(name):
    print("Hello", name)
    return name.upper()
\end{dhkcode}

원본 \zsh{greet} 함수 객체가 생성된다.

여기서 중요한 점은 함수 객체가 생성되었다고 함수 본문이 실행되는 것은 아니라는 사실이다.

따라서 아직 \zsh{"Hello Kim"}은 출력되지 않는다.

\begin{dhknote}
\begin{dhksmry}[centered]
2. 데코레이터 적용
\end{dhksmry}
\end{dhknote}

생성된 함수 객체가 \zsh{register}에 전달된다.

\begin{dhkcode}
register(original_greet)
\end{dhkcode}

\zsh{register} 내부에서는 다음 문장이 실행된다.

\begin{dhkcode}
registry[func.__name__] = func
\end{dhkcode}

\zsh{func.__name__}은 \zsh{"greet"}이므로 레지스트리에 다음 관계가 형성된다.

\begin{dhkdir}[in=6]
registry
> "greet" -> 원본 greet 함수 객체
\end{dhkdir}

그리고 \zsh{register}는 같은 함수 객체를 반환한다.

\begin{dhknote}
\begin{dhksmry}[centered]
3. 최종 이름 바인딩
\end{dhksmry}
\end{dhknote}

데코레이터의 반환값이 모듈 네임스페이스의 \zsh{greet} 이름에 바인딩된다.

따라서 다음 두 표현식은 동일한 함수 객체를 참조한다.

\begin{dhkcode}
greet
registry["greet"]
\end{dhkcode}

이것이 다음 결과가 \zsh{True}인 이유다.

\begin{dhkcode}
registry["greet"] is greet
\end{dhkcode}

\begin{dhknote}
\begin{dhksmry}[centered]
4. 레지스트리를 통한 호출
\end{dhksmry}
\end{dhknote}

마지막으로 다음 표현식이 실행된다.

\begin{dhkcode}
registry["greet"]("Kim")
\end{dhkcode}

딕셔너리 조회 결과는 함수 객체다.

따라서 일반적인 함수 호출 규칙이 적용된다.

이때 비로소 원본 함수 본문이 실행되고 \zsh{"Hello Kim"}이 출력된다.

\begin{dhknote}
\begin{dhksmry}[centered]
여기서 발생한 상태 변화
\end{dhksmry}
\end{dhknote}

이 예제에서는 원본 함수 객체를 수정한 것이 아니다.

변경된 것은 외부 딕셔너리의 상태다.

\begin{dhktable}
|대상	|실행 전	|실행 후|
|:---|:---|:---|
|원본 greet 함수 객체|	생성 전|	생성됨|
|registry 딕셔너리|	빈 딕셔너리|	"greet" → 함수 객체|
|모듈 이름 greet|	정의 전|	원본 함수 객체 참조|
|원본 함수 본문|	미실행|	레지스트리에서 호출할 때 실행|
\end{dhktable}

이것이 wrapper 데코레이터와 등록형 데코레이터를 구별하는 첫 번째 기준이다.

\subsubsection{등록형 데코레이터에도 인수를 전달할 수 있다}

\zsh{A.9}에서는 데코레이터 팩토리를 학습했다.

등록형 데코레이터도 같은 구조로 확장할 수 있다.

\begin{dhkcode}
registry = {}


def register(name):
    print("factory:", name)

    def decorator(func):
        print("register:", func.__name__)

        registry[name] = func

        return func

    return decorator


@register("hello")
def greet(name):
    return f"Hello, {name}"


print("definition complete")

print(registry["hello"] is greet)
print(registry["hello"]("Kim"))
\end{dhkcode}

예상 출력:

\begin{dhkzsh}
factory: hello
register: greet
definition complete
True
Hello, Kim
\end{dhkzsh}

이제 데코레이터의 실행 과정을 세 단계로 구분할 수 있다.

\begin{dhksmry}
데코레이터 팩토리 \ra 등록 \ra 호출

1. 팩토리 호출
   register("hello")가 실행되어 "hello"를 클로저로 유지하는
   decorator 함수 객체를 반환한다.

2. 등록 수행
   원본 greet 함수 객체가 생성된 뒤 decorator(greet)이 호출되고,
   registry["hello"]에 함수 객체가 저장된다.

3. 실제 함수 호출
   이후 registry["hello"]("Kim")을 실행할 때 비로소 원본 함수
   본문이 실행된다.
\end{dhksmry}

여기서 \zsh{register("hello")}의 반환값은 데코레이터 함수 객체다.

반면 그 데코레이터의 반환값은 원본 \zsh{greet} 함수 객체다.

따라서 다음 두 \zsh{return}은 의미가 다르다.

\begin{dhkcode}
return decorator
\end{dhkcode}

이것은 설정이 완료된 데코레이터 함수 객체를 반환하는 단계다.

\begin{dhkcode}
return func
\end{dhkcode}

이것은 등록된 원본 함수 객체를 반환하는 단계다.

이 구분을 이해했다면 FastAPI의 \zsh{@app.get(...)} 구조도 같은 원리로 분석할 수 있다.

\subsubsection{FastAPI의 \zsh{@app.get(...)}은 어떤 구조일까?}

이제 실제 FastAPI를 살펴보자.

\begin{dhkcode}
from fastapi import FastAPI

app = FastAPI()


@app.get("/items/{item_id}")
def get_item(item_id: int):
    return {"item_id": item_id}
\end{dhkcode}

여기서 가장 먼저 구별해야 할 객체는 세 가지다.

\begin{dhktable}
|표현식|	객체의 역할|
|:---|:---|
|app|	FastAPI 애플리케이션 인스턴스|
|app.get|	app에 바인딩된 메서드|
|app.get("/items/{item_id}")|	경로 등록을 수행할 데코레이터 함수|
\end{dhktable}

\zsh{A.11}에서 학습한 메서드 바인딩을 연결해 보자.

\begin{dhkcode}
print(app.get.__self__ is app)
\end{dhkcode}

예상 출력:

\begin{dhkzsh}
True
\end{dhkzsh}

\zsh{app.get}은 FastAPI 클래스에 정의된 메서드가 \zsh{app} 인스턴스에 바인딩된 결과다.

따라서 다음 표현식은

\begin{dhkcode}
app.get("/items/{item_id}")
\end{dhkcode}

FastAPI 인스턴스의 \zsh{get} 메서드를 호출하는 것이다.

그 호출 결과로 데코레이터 함수 객체가 반환된다.

여기서 \zsh{get} 메서드 호출과 엔드포인트 함수 호출은 완전히 다른 사건이다.

\zsh{app.get(...)}을 실행한다고 \zsh{get_item} 함수의 본문이 실행되는 것은 아니다.

\begin{dhknote}
\begin{dhksmry}[centered]
실제 변환식
\end{dhksmry}
\end{dhknote}

다음 코드는

\begin{dhkcode}
@app.get("/items/{item_id}")
def get_item(item_id: int):
    return {"item_id": item_id}
\end{dhkcode}

개념적으로 다음과 대응한다.

\begin{dhkcode}
decorator = app.get("/items/{item_id}")


def get_item(item_id: int):
    return {"item_id": item_id}


get_item = decorator(get_item)
\end{dhkcode}

FastAPI의 일반적인 경로 등록 데코레이터는 전달받은 엔드포인트 함수 객체를 라우터에 등록한 뒤 그 함수 객체를 그대로 반환한다.

따라서 데코레이션 과정에서 별도의 사용자 정의 wrapper가 반드시 생성되는 것은 아니다.

다만 FastAPI는 등록된 엔드포인트를 HTTP 요청으로 실행하기 위한 내부 요청 처리 함수를 별도로 구성한다.

등록 데코레이터가 원본 함수를 반환한다는 사실과 프레임워크 내부에 요청 처리 계층이 존재한다는 사실은 양립한다.

\subsubsection{등록 시점과 HTTP 요청 처리 시점}

FastAPI의 실행 과정을 두 시점으로 나누어 보자.

\begin{dhknote}
\begin{dhksmry}[centered]
1. 등록 시점 - Python 모듈 실행
\end{dhksmry}
\end{dhknote}

\begin{dhkgraph}[uniform]
node n1 (1,1) :: app = FastAPI()
group g1(n1) :: 애플리케이션 인스턴스 생성
node n2 (2,1) :: app.get("/items/{item_id}")
group g2(n2) :: 등록 데코레이터 함수 생성
node n3 (3,1) :: def get_item(...)
group g3(n3) :: 엔드포인트 함수 객체 생성
node n4 (4,1) :: 데코레이터 적용 -> 라우터 등록
group g4(n4) :: 경로, HTTP 메서드, 엔드포인트 참조 저장

edge g1 -> g2
edge g2 -> g3
edge g3 -> g4
\end{dhkgraph}

\begin{dhknote}
\begin{dhksmry}[centered]
2. 요청 처리 시점 - HTTP GET 요청 도착
\end{dhksmry}
\end{dhknote}

\begin{dhkgraph}[uniform]
node n1 (1,1) :: 요청 경로와 HTTP 메서드에 맞는 라우트 탐색
group g1(n1) :: GET /items/7 
node n2 (2,1) :: 경로의 "7"을 정수 7로 변환
group g2(n2) :: 매개변수 추출·검증·의존성 해결 
node n3 (3,1) :: 등록된 함수 객체를 실제로 호출
group g3(n3) :: get_item(item_id=7)
node n4 (4,1) :: 반환값을 응답 데이터로 직렬화
group g4(n4) :: HTTP 응답 생성

edge g1 -> g2
edge g2 -> g3
edge g3 -> g4
\end{dhkgraph}

여기서 중요한 것은 등록과 호출 사이에 시간적 간격이 존재한다는 사실이다.

보통 Python 모듈이 실행되어 경로가 등록된 뒤, 나중에 HTTP 요청이 도착했을 때 엔드포인트가 호출된다.

같은 애플리케이션 프로세스에서 요청이 여러 번 발생한다고 해서 매번 함수 정의문이나 \zsh{@app.get(...)}이 다시 실행되는 것은 아니다.

\subsubsection{실제 FastAPI의 등록 객체를 확인하기}

이번에는 FastAPI 내부에 어떤 객체가 저장되는지 확인해 보자.

다음 예제는 \zsh{fastapi}와 \zsh{httpx}가 설치된 환경에서 독립적으로 실행할 수 있다.

\zsh{TestClient}는 별도의 웹 서버를 실행하지 않고 애플리케이션에 테스트 요청을 전달하는 도구다.

\begin{dhkcode}
from fastapi import FastAPI
from fastapi.routing import APIRoute
from fastapi.testclient import TestClient

app = FastAPI()

@app.get("/items/{item_id}")
def get_item(item_id: int):
    print("endpoint called:", item_id)
    return {"item_id": item_id}

route = next(
    route
    for route in app.routes
    if isinstance(route, APIRoute)
    and route.path == "/items/{item_id}"
)

print(route.endpoint is get_item)

client = TestClient(app)

response = client.get("/items/7")

print(response.status_code)
print(response.json())

invalid = client.get("/items/abc")

print(invalid.status_code)
\end{dhkcode}

예상 출력:

\begin{dhkzsh}
True
endpoint called: 7
200
{'item_id': 7}
422
\end{dhkzsh}

이 예제에는 서로 다른 두 종류의 함수 호출이 등장한다.

\begin{dhknote}
\begin{dhksmry}[centered]
1. 등록된 엔드포인트의 동일성
\end{dhksmry}
\end{dhknote}

\begin{dhkcode}
route.endpoint is get_item
\end{dhkcode}

결과는 \zsh{True}다.

\zsh{route}는 FastAPI의 \zsh{APIRoute} 객체이고, \zsh{route.endpoint}는 등록된 엔드포인트 함수 객체를 참조한다.

현재 예제에서는 \zsh{@app.get(...)}이 원본 함수 객체를 반환했으므로 모듈 이름 \zsh{get_item}도 같은 객체를 참조한다.

따라서 다음 구조가 형성된다.

\begin{dhkdiagram}[in=6]
node n1 (1,1) :: 모듈 네임스페이스
node n2 (2,1) :: get_item
node n3 (3,2) :: 함수 객체
node n4 (4,1) :: FastAPI 애플리케이션
node n5 (5,1) :: APIRoute 객체
node n6 (6,1) :: endpoint

edge n1 -- n2
edge n2 -> n3 :: HV
edge n4 -- n5
edge n5 -- n6
edge n6 -> n3 :: HV
\end{dhkdiagram}

여기서 함수 객체가 복사된 것은 아니다.

서로 다른 두 위치에서 동일한 함수 객체를 참조하는 것이다.

\begin{dhknote}
\begin{dhksmry}[centered]
2. 실제 HTTP 요청
\end{dhksmry}
\end{dhknote}

\begin{dhkcode}
response = client.get("/items/7")
\end{dhkcode}

이 코드는 등록된 함수 객체를 직접 호출하는 코드가 아니다.

HTTP 요청을 모사하여 FastAPI 애플리케이션의 요청 처리 경로로 전달한다.

FastAPI는 요청 경로에서 \zsh{"7"}을 추출하고, 엔드포인트의 타입 힌트를 근거로 정수 7을 준비한다.

그다음 엔드포인트 함수를 호출한다.

이때 비로소 다음 출력이 발생한다.

\begin{dhkzsh}
endpoint called: 7
\end{dhkzsh}

반면 다음 요청은

\begin{dhkcode}
client.get("/items/abc")
\end{dhkcode}

\zsh{"abc"}를 정수로 변환할 수 없으므로 요청 매개변수 검증 단계에서 실패한다.

따라서 HTTP 상태 코드 422가 반환되고, 해당 요청에서는 \zsh{get_item} 함수 본문이 실행되지 않는다.

이것이 함수 객체의 등록과 함수 본ㅁ누의 실행이 서로 다른 시점에 발생한다는 사실을 보여주는 예제다.

\subsubsection{직접 함수 호출과 HTTP 요청은 같지 않다}

앞의 예제에서 다음 두 표현식을 비교해 보자.

\begin{dhkcode}
get_item("7")

client.get("/items/7")
\end{dhkcode}

첫 번째는 일반적인 Python 함수 호출이다.

두 번째는 FastAPI의 HTTP 요청 처리 과정이다.

따라서 동작이 다르다.

\begin{dhkcode}
print(get_item("7"))
\end{dhkcode}

예상 출력:

\begin{dhkzsh}
endpoint called: 7
{'item_id': '7'}
\end{dhkzsh}

여기서는 \zsh{item_id}에 문자열 \zsh{"7"}이 그대로 전달된다.

Python의 타입 힌트 \zsh{item_id: int} 자체는 일반적인 함수 호출에서 자동 형 변환을 수행하지 않는다.

반면 HTTP 요청에서는 FastAPI가 엔드포인트 시그니처를 분석하여 입력값을 검증하고 변환한다.

\begin{dhktable}
|구분	|직접 함수 호출|	FastAPI HTTP 요청|
|:---|:---|:---|
|함수 호출 주체|	Python 코드|	FastAPI 요청 처리 계층|
|경로 매개변수 추출|	없음|	있음|
|타입 힌트 기반 입력 검증|	자동 수행되지 않음|	수행|
|의존성 주입|	자동 수행되지 않음|	수행|
|HTTP 응답 생성|	없음	|수행|
|원본 함수 본문|	직접 실행|	검증·의존성 처리 후 실행|
\end{dhktable}

여기서 중요한 점은 FastAPI가 Python 함수 호출 규칙을 변경한 것이 아니라는 사실이다.

FastAPI는 함수를 호출하기 전에 HTTP 요청을 분석하고, 필요한 인수들을 구성하는 별도의 계층을 제공하는 것이다.

\subsubsection{FastAPI는 함수 시그니처를 어떻게 활용할까?}

이제 \zsh{A.8}에서 학습한 \zsh{inspect.signature()}와 연결해 보자.

FastAPI는 등록된 함수 매개변수, 타입 힌트, 기본값 등을 분석하여 HTTP 요청에서 어떤 값을 추출하고 어떤 의존성을 해결해야 하는지 결정한다.

다음 예제를 보자.

\begin{dhkcode}
import inspect
from fastapi import FastAPI, Depends
from fastapi.testclient import TestClient

app = FastAPI()

def get_multiplier() -> int:
    return 3

@app.get("/scores/{score}")
def calculate_score(
    score: int,
    offset: int = 0,
    multiplier: int = Depends(get_multiplier),
):
    return {
        "result": (score + offset) * multiplier
    }

signature = inspect.signature(calculate_score)

print(signature.parameters["score"].annotation)
print(signature.parameters["offset"].default)

print(
    type(
        signature.parameters["multiplier"].default
    ).__name__
)

client = TestClient(app)

response = client.get("/scores/4?offset=2")

print(response.json())
\end{dhkcode}

예상 출력:

\begin{dhkzsh}
<class 'int'>
0
Depends
{'result': 18}
\end{dhkzsh}

이 코드에서 세 매개변수는 서로 다른 방식으로 처리된다.

\begin{dhktable}
|매개변수|	선언|	FastAPI에서의 역할|
|:---|:---|:---|
|score|	score: int|	URL 경로에서 추출|
|offset|	offset: int = 0|	쿼리 매개변수에서 추출, 생략 시 기본값 사용|
|multiplier|	multiplier: int = Depends(//*get_multiplier)|	의존성 함수의 반환값을 전달|
\end{dhktable}

요청은 다음과 같다.

\begin{dhkzsh}
GET /scores/4?offset=2
\end{dhkzsh}

FastAPI는 등록된 경로와 함수 시그니처를 근거로 인수를 준비한다.

\begin{dhkzsh}
score = 4
offset = 2
multiplier = 3
\end{dhkzsh}

여기서 \zsh{multiplier}의 값은 \zsh{get_multiplier()}를 호출하여 얻는다.

따라서 엔드포인트 호출은 개념적으로 다음과 같다.

\begin{dhkcode}
calculate_score(
    score=4,
    offset=2,
    multiplier=3,
)
\end{dhkcode}

그 결과는 다음과 같다.

\[
(4+2)\times3=18
\]

\begin{dhknote}
\begin{dhksmry}[centered]
Depends(...)는 실제 의존성의 반환값이 아니다
\end{dhksmry}
\end{dhknote}

여기서 중요한 점이 있다.

다음 선언을 보자.

\begin{dhkcode}
multiplier: int = Depends(get_multiplier)
\end{dhkcode}

일반적인 Python 함수 정의 관점에서는 \zsh{Depends(get_multiplier)}의 호출 결과인 \zsh{Depends} 객체가 기본값으로 저장된다.

이 객체 자체가 정수 3인 것은 아니다.

FastAPI는 이 기본값을 의존성 해결을 위한 선언 정보로 해석한다.

따라서 다음처럼 직접 호출하면

\begin{dhkcode}
calculate_score(4, offset=2)
\end{dhkcode}

FastAPI의 의존성 주입 과정이 실행되지 않는다.

\zsh{multiplier}에는 기본값인 \zsh{Depends} 객체가 바인딩되므로, 이 예제에서는 정수와 \zsh{Depends} 객체를 곱하려는 연산때문에 \zsh{TypeError}가 발생한다.

반면 다음 직접 호출은 가능하다.

\begin{dhkcode}
calculate_score(
    4,
    offset=2,
    multiplier=3,
)
\end{dhkcode}

이처럼 FastAPI의 의존성 주입은 Python의 함수 호출 규칙 자체가 아니라 프레임워크가 함수 호출 전에 수행하는 별도의 인수 구성 과정이다.

\subsubsection{\zsh{@app.get}과 \zsh{@trace}의 적용 순서를 바꾸면?}

이번 학습에서 가장 중요한 실전 문제다.

\zsh{A.10}에서는 여러 데코레이터의 적용 순서가 최종 객체 구조를 결정한다는 사실을 학습했다.

FastAPI에서도 같은 원리가 적용된다.

먼저 일반적인 로깅 데코레이터를 정의해 보자.

\begin{dhkcode}
from functools import wraps

from fastapi import FastAPI
from fastapi.routing import APIRoute
from fastapi.testclient import TestClient


app = FastAPI()


def trace(func):

    @wraps(func)
    def wrapper(*args, **kwargs):
        print("TRACE:", func.__name__)
        return func(*args, **kwargs)

    return wrapper
\end{dhkcode}

이제 적용 순서를 다르게 하여 두 엔드포인트를 등록해 보자.

\begin{dhkcode}
@app.get("/logged")
@trace
def logged():
    return {"name": "logged"}


@trace
@app.get("/unlogged")
def unlogged():
    return {"name": "unlogged"}
\end{dhkcode}

겉으로는 두 함수 모두 \zsh{@app.get}과 \zsh{@trace}를 사용했다.

하지만 FastAPI에 등록되는 함수 객체는 다르다.

\begin{dhknote}
\begin{dhksmry}[centered]
첫 번째 - @app.get이 바깥쪽
\end{dhksmry}
\end{dhknote}

\begin{dhkcode}
@app.get("/logged")
@trace
def logged():
    ...
\end{dhkcode}

데코레이터의 적용 순서는 아래에서 위다.

따라서 개념적으로 다음과 같다.

\begin{dhkcode}
logged = app.get("/logged")(trace(original_logged))
\end{dhkcode}

실행 과정을 추적해 보자.

먼저 원본 \zsh{logged} 함수 객체가 생성된다.

그다음 \zsh{trace}가 원본 함수를 전달받아 wrapper를 반환한다.

이제 \zsh{app.get("/logged")}가 반환한 등록 데코레이터가 wrapper 함수 객체를 엔드포인트로 등록한다.

마지막으로 wrapper 함수 객체가 모듈 이름 \zsh{logged}에 바인딩된다.

따라서 다음 관계가 성립한다.

\begin{dhkdiagram}[in=2]
node n1 (1,1) :: 모듈 이름 logged
node n2 (3,1) :: trace wrapper
node n3 (3,2) :: FastAPI route.endpoint
node n4 (5,1) :: 원본 logged 함수

edge n1 -> n2
edge n3 -> n2
edge n2 -> n4 :: 클로저의 func
\end{dhkdiagram}

HTTP 요청이 들어오면 FastAPI가 wrapper를 호출한다.

Wrapper는 \zsh{"TRACE: logged"}를 출력한 뒤 원본 함수를 호출한다.

\begin{dhknote}
\begin{dhksmry}[centered]
두 번째 - @trace가 바깥쪽
\end{dhksmry}
\end{dhknote}

\begin{dhkcode}
@trace
@app.get("/unlogged")
def unlogged():
    ...
\end{dhkcode}

이번에는 다음과 대응한다.

\begin{dhkcode}
unlogged = trace(
    app.get("/unlogged")(original_unlogged)
)
\end{dhkcode}

먼저 원본 함수 객체가 생성된다.

그다음 \zsh{@app.get}이 원본 함수 객체를 FastAPI에 등록하고 그대로 반환한다.

이후 \zsh{trace}가 그 원본 함수 객체를 감싸 새로운 wrapper를 반환한다.

마지막으로 모듈 이름 \zsh{unlogged}가 wrapper에 바인딩된다.

따라서 구조는 다음과 같다.

\begin{dhkdiagram}[in=2]
node n1 (1,1) :: 모듈 이름 unlogged
node n2 (3,1) :: trace wrapper
node n3 (5,2) :: FastAPI route.endpoint
node n4 (5,1) :: 원본 unlogged 함수

edge n1 -> n2
edge n3 -> n4
edge n2 -> n4 :: 클로저의 func
\end{dhkdiagram}

이번에는 FastAPI가 원본 함수 객체를 등록했다.

따라서 HTTP 요청을 처리할 때 wrapper를 거치지 않는다.

하지만 Python 코드에서 \zsh{unlogged()}를 직접 호출하면 wrapper가 실행된다.

\begin{dhknote}
\begin{dhksmry}[centered]
실제 객체 동일성 검증
\end{dhksmry}
\end{dhknote}

앞의 코드에 이어 다음을 실행해 보자.

\begin{dhkcode}
def endpoint_for(path):
    return next(
        route.endpoint
        for route in app.routes
        if isinstance(route, APIRoute)
        and route.path == path
    )


print(endpoint_for("/logged") is logged)

print(endpoint_for("/unlogged") is unlogged)

print(
    endpoint_for("/unlogged")
    is unlogged.__wrapped__
)


client = TestClient(app)

client.get("/logged")
client.get("/unlogged")

unlogged()
\end{dhkcode}

예상 출력:

\begin{dhkzsh}
True
False
True
TRACE: logged
TRACE: unlogged
\end{dhkzsh}

출력에서 \zsh{"TRACE: logged"}는 HTTP 요청에 의해 발생한다.

반면 \zsh{"TRACE: unlogged"}는 마지막의 직접 함수 호출에 의해 발생한다.

\zsh{client.get("/unlogged")} 자체는 로그를 출력하지 않는다.

이 결과는 단순한 실행 순서 차이가 아니다.

FastAPI가 참조하는 함수 객체와 모듈 이름이 참조하는 함수 객체가 서로 달라진 결과다.

\begin{dhknote}
\begin{dhksmry}[centered]
두 순서의 비교
\end{dhksmry}
\end{dhknote}

\begin{dhktable}
|구분|	@app.get 바깥쪽|	@trace 바깥쪽|
|:---|:---|:---|
|등록되는 함수 객체|	trace wrapper|	원본 함수|
|최종 함수 이름|	trace wrapper|	trace wrapper|
|HTTP 요청에서 trace 실행|	실행됨|	실행되지 않음|
|직접 함수 호출에서 trace 실행|	실행됨	|실행됨|
|route.endpoint is 함수이름|	True|	False|
\end{dhktable}

이것이 \zsh{A.10}에서 학습한 데코레이터 적용 순서가 실제 웹 프레임워크의 동작에 영향을 주는 대표적인 사례다.

\subsubsection{\zsh{functools.wraps}가 FastAPI에서 중요한 이유}

앞의 \zsh{trace} 구현에는 다음 코드가 있었다.

\begin{dhkcode}
@wraps(func)
def wrapper(*args, **kwargs):
    ...
\end{dhkcode}

왜 \zsh{@wraps(func)}를 사용했을까?

\zsh{A.8}에서 학습했듯 \zsh{functools.wraps}는 wrapper의 메타데이터를 설정하고 \zsh{__wrapped__}를 통해 원본 함수 객체를 추적할 수 있도록 한다.

FastAPI는 엔드포인트의 시그니처를 분석하므로 이 정보가 중요하다.

다음 코드를 살펴보자.

\begin{dhkcode}
import inspect
from functools import wraps


def trace(func):

    @wraps(func)
    def wrapper(*args, **kwargs):
        return func(*args, **kwargs)

    return wrapper


@trace
def calculate(x: int, y: int = 1) -> int:
    return x + y


print(inspect.signature(calculate))

print(
    inspect.signature(
        calculate,
        follow_wrapped=False
    )
)
\end{dhkcode}

예상 출력:

\begin{dhkzsh}
(x: int, y: int = 1) -> int
(*args, **kwargs)
\end{dhkzsh}

두 결과가 다른 이유는 다음과 같다.

\zsh{inspect.signature(calculate)}는 기본적으로 \zsh{__wrapped__} 체인을 따라 원본 함수의 시그니처를 분석한다.

반면 \zsh{follow_wrapped=False}를 지정하면 실제 wrapper 함수 객체의 시그니처를 확인한다.

이 차이가 FastAPI에서 중요하다.

\begin{dhknote}
\begin{dhksmry}[centered]
@wraps가 없으면 어떤 문제가 발생할까?
\end{dhksmry}
\end{dhknote}

다음 코드를 살펴보자.

\begin{dhkcode}
from fastapi import FastAPI
from fastapi.testclient import TestClient


app = FastAPI()


def bad_trace(func):

    def wrapper(*args, **kwargs):
        return func(*args, **kwargs)

    return wrapper


@app.get("/broken/{item_id}")
@bad_trace
def broken(item_id: int):
    return {"item_id": item_id}


client = TestClient(app)

response = client.get("/broken/5")

print(response.status_code)
\end{dhkcode}

예상 출력:

\begin{dhkzsh}
422
\end{dhkzsh}

왜 정상적으로 실행되지 않을까?

FastAPI에 등록된 함수 객체는 \zsh{bad_trace}가 반환한 wrapper다.

그런데 이 wrapper의 시그니처는 다음과 같다.

\begin{dhkcode}
(*args, **kwargs)
\end{dhkcode}

\zsh{__wrapped__} 참조도 설정되어 있지 않다.

따라서 FastAPI는 원본 함수의 \zsh{item_id: int} 매개변수를 정상적으로 추적하지 못하고, wrapper의 \zsh{args}, \zsh{kwargs}를 요청 매개변수로 분석할 수 있다.

이 예제에서는 필요한 쿼리 매개변수가 제공되지 않았다고 판단하여 422 응답이 발생한다.

\zsh{@wraps(func)}를 사용하면 원본 함수의 시그니처를 추적할 수 있으므로 이런 문제를 방지하는 데 도움이 된다.

하지만 한 가지 주의해야 한다.

\zsh{@wraps}는 메타데이터를 설정하는 기능이지, wrapper의 실제 호출 동작을 원본 함수와 동일하게 만들어 주는 기능은 아니다.

예를 들어 원본 함수의 시그니처를 노출하더라도 wrapper가 전달받은 인수를 올바르게 처리하지 않으면 실행 시 오류가 발생할 수 있다.

따라서 FastAPI용 데코레이터는 메타데이터와 실제 인수 전달 규칙을 모두 고려해야 한다.

\subsubsection{비동기 엔드포인트와 데코레이터}

이제 FastAPI에서 특히 중요한 문제를 살펴보자.

FastAPI는 다음 두 종류의 엔드포인트를 지원한다.

\begin{dhkcode}
def sync_endpoint():
    ...

async def async_endpoint():
    ...
\end{dhkcode}

두 함수는 호출 결과와 실행 방식이 다르다.

일반적인 \zsh{def} 함수는 호출하면 본문을 실행하고 반환값을 돌려준다.

반면 \zsh{async def}로 정의한 코루틴 함수(Coroutine Function)를 호출하면 코루틴 객체가 반환된다.

함수 본문의 실행을 진행하려면 해당 코루틴을 \zsh{await}하는 등의 과정이 필요하다.

이 차이는 데코레이터에도 영향을 준다.

\begin{dhknote}
\begin{dhksmry}[centered]
비동기 함수를 동기 wrapper로 감싸면?
\end{dhksmry}
\end{dhknote}

다음 코드를 살펴보자.

\begin{dhkcode}
from functools import wraps
import inspect


def trace(func):

    @wraps(func)
    def wrapper(*args, **kwargs):
        print("TRACE")
        return func(*args, **kwargs)

    return wrapper


@trace
async def calculate(x):
    return x * 2


print(inspect.iscoroutinefunction(calculate))
\end{dhkcode}

예상 출력:

\begin{dhkzsh}
False
\end{dhkzsh}

왜 \zsh{False}일까?

원본 \zsh{calculate}는 코루틴 함수다.

하지만 데코레이션 이후 \zsh{calculate} 이름은 일반적인 \zsh{def wrapper} 함수 객체를 참조한다.

\zsh{@wraps(func)}가 \zsh{__wrapped__}를 설정했다고 해서 wrapper 자체가 \zsh{async def} 함수로 변환되는 것은 아니다.

이제 다음 표현식을 생각해 보자.

\begin{dhkcode}
result = calculate(10)
\end{dhkcode}

일반 wrapper가 실행되면서 원본 코루틴 함수를 호출한다.

원본 함수 호출 결과는 코루틴 객체다.

Wrapper는 그 코루틴 객체를 그대로 반환한다.

따라서 \zsh{result}는 정수 20이 아니라 코루틴 객체다.

이 코루틴을 실행하여 결과를 얻으려면 별도로 \zsh{await} 해야 한다.

\begin{dhknote}
\begin{dhksmry}[centered]
FastAPI에서는 왜 문제가 될까?
\end{dhksmry}
\end{dhknote}

FastAPI는 엔드포인트가 동기 함수인지 비동기 함수인지에 따라 실행 방식을 달리한다.

일반적인 동기 엔드포인트는 스레드풀에서 실행하고, 비동기 엔드포인트는 비동기 실행 경로에서 \zsh{await}한다.

그런데 코루틴 함수를 동기 wrapper로 감싸면 FastAPI가 최종 등록된 wrapper를 동기 함수로 판단할 수 있다.

이 경우 wrapper 호출 결과로 코루틴 객체가 반환되지만, 그 객체가 정상적으로 \zsh{await} 되지 않은 채 응답 직렬화 단계에 전달되어 오류가 발생할 수 있다.

따라서 \zsh{@wraps}만 사용했다고 해서 비동기 엔드포인트에 안전한 데코레이터가 되는 것은 아니다.

\begin{dhknote}
\begin{dhksmry}[centered]
동기·비동기 함수를 구별하는 데코레이터
\end{dhksmry}
\end{dhknote}

다음처럼 두 종류의 wrapper를 구별할 수 있다.

\begin{dhkcode}
import inspect
from functools import wraps

from fastapi import FastAPI
from fastapi.testclient import TestClient


def trace(func):

    if inspect.iscoroutinefunction(func):

        @wraps(func)
        async def async_wrapper(*args, **kwargs):
            print("TRACE async:", func.__name__)

            result = await func(*args, **kwargs)

            return result

        return async_wrapper

    @wraps(func)
    def sync_wrapper(*args, **kwargs):
        print("TRACE sync:", func.__name__)

        result = func(*args, **kwargs)

        return result

    return sync_wrapper


app = FastAPI()


@app.get("/sync")
@trace
def sync_endpoint():
    return {"mode": "sync"}


@app.get("/async")
@trace
async def async_endpoint():
    return {"mode": "async"}


client = TestClient(app)

print(client.get("/sync").json())
print(client.get("/async").json())
\end{dhkcode}

여기서 두 wrapper는 서로 다른 함수 객체다.

동기 엔드포인트에 적용된 wrapper는 일반 함수 객체이고, 비도익 엔드포인트에 적용된 wrapper는 코루틴 함수 객체다.

따라서 FastAPI는 각 엔드포인트의 실행 방식을 적절하게 구별할 수 있다.

다만 이 구현은 직접 전달받은 일반적인 동기 함수와 코루틴 함수를 처리하는 예제다.

이미 다른 데코레이터로 감싸진 함수, 비동기 제네레이터, 복잡한 호출 가능 객체까지 모두 처리하는 범용 구현은 아니다.

여기서 핵심은 다음과 같다.

\begin{dhksmry}[centered]
함수 객체의 메타데이터, 실제 호출 결과, 

동기·비동기 실행 방식은 각각 구별해서 분석해야 한다.
\end{dhksmry}

\subsubsection{\zsh{APIRouter} - 등록 정보를 모듈 단위로 분리하기}

FastAPI 애플리케이션이 커지면 모든 엔드포인트를 하나의 파일에 정의하기 어려워진다.

이때 \zsh{APIRouter}를 사용할 수 있다.

\begin{dhkcode}
from fastapi import FastAPI, APIRouter
from fastapi.testclient import TestClient


router = APIRouter(prefix="/v1")


@router.get("/items/{item_id}")
def get_item(item_id: int):
    return {"item_id": item_id}


app = FastAPI()

app.include_router(router)


client = TestClient(app)

response = client.get("/v1/items/10")

print(response.status_code)
print(response.json())
\end{dhkcode}

예상 출력:

\begin{dhkzsh}
200
{'item_id': 10}
\end{dhkzsh}

이 구조에서 다음 두 과정은 구별해야 한다.

\begin{dhknote}
\begin{dhksmry}[centered]
첫째, @router.get(...)에 의한 등록
\end{dhksmry}
\end{dhknote}

\begin{dhkcode}
@router.get("/items/{item_id}")
\end{dhkcode}

이 데코레이터는 엔드포인트 함수를 \zsh{router}에 등록한다.

\begin{dhknote}
\begin{dhksmry}[centered]
둘째, app.include_router(router)에 의한 애플리케이션 구성
\end{dhksmry}
\end{dhknote}

\begin{dhkcode}
app.include_router(router)
\end{dhkcode}

이 호출은 해당 라우터에 정의된 경로들을 FastAPI 애플리케이션에 포함시킨다.

따라서 전체 경로는 다음과 같이 구성된다.

\begin{dhksmry}[centered]
APIRouter prefix "/v1"
+
등록된 경로 "/items/{item_id}"
\da
"/v1/items/{item_id}"
\end{dhksmry}

이 구조를 객체 참조 관점에서 생각해 보자.

\zsh{router}와 \zsh{app}은 서로 다른 객체다.

\zsh{@router.get(...)}은 먼저 \zsh{router}에 경로 정보를 등록한다.

그다음 \zsh{app.include_router(router)}가 라우터의 등록 정보를 애플리케이션에 반영한다.

FastAPI의 일반적인 \zsh{include_router} 구현에서는 포함 시점에 라우터의 경로들을 애플리케이션 쪽 라우터에 등록하는 방식으로 처리한다.

따라서 이미 \zsh{include_router}를 실행한 뒤 원래 \zsh{router}에 새로운 경로를 추가했다고 해서 그 경로가 기존 애플리케이션에 자동 반영된다고 가정하면 안 된다.

이것이 모듈을 나누어 FastAPI 애플리케이션을 구성할 때 등록 시점을 고려해야 하는 이유다.

\subsubsection{하나의 함수 객체를 여러 경로에 등록할 수 있을까?}

가능하다.

다음 코드를 살펴보자.

\begin{dhkcode}
from fastapi import FastAPI
from fastapi.routing import APIRoute


app = FastAPI()


@app.get("/first")
@app.get("/second")
def greet():
    return {"message": "Hello"}


routes = [
    route
    for route in app.routes
    if isinstance(route, APIRoute)
]

for route in routes:
    print(route.path, route.endpoint is greet)
\end{dhkcode}

예상 출력:

\begin{dhkzsh}
/second True
/first True
\end{dhkzsh}

왜 \zsh{/second}가 먼저 등록될까?

\zsh{A.10}에서 학습한 데코레이터 적용 순서 때문이다.

데코레이ㅓㅌ 표현식은 위에서 아래로 평가되지만, 함수 객체에 대한 적용은 아래에서 위로 진행된다.

따라서 개념적으로 다음과 같다.

\begin{dhkcode}
greet = app.get("/first")(
    app.get("/second")(greet)
)
\end{dhkcode}

먼저 \zsh{/second}가 등록되고, 그다음 \zsh{/first}가 등록된다.

그리고 두 등록 데코레이터 모두 원본 함수 객체를 그대로 반환한다.

따라서 두 라우트의 \zsh{endpoint}는 동일한 함수 객체를 참조한다.

\begin{dhkdiagram}[in=6]
node n1 (1,1) :: GET /first
node n2 (3,1) :: GET /second
node n3 (2,2) :: greet 함수 객체

edge n1 -> n3 :: HV
edge n2 -> n3 :: HV
\end{dhkdiagram}

이 구조는 하나의 함수를 여러 URL에 연결할 때 활용할 수 있다.

다만 라우트 등록 순서는 실제 요청 매칭에도 영향을 줄 수 있다.

예를 들어 고정 경로와 동적 경로가 서로 겹치는 경우에는 어떤 경로를 먼저 등록했는지 고려해야 한다.

\subsubsection{등록 이후 함수 이름을 재바인딩하면 어떻게 될까?}

등록형 데코레이터의 가장 중요한 특징 중 하나는 외부 레지스트리가 함수 객체에 대한 참조를 보유한다는 점이다.

다음 코드를 보자.

\begin{dhkcode}
from fastapi import FastAPI
from fastapi.routing import APIRoute


app = FastAPI()


@app.get("/hello")
def greet():
    return {"message": "Original"}


route = next(
    route
    for route in app.routes
    if isinstance(route, APIRoute)
    and route.path == "/hello"
)


registered = route.endpoint


def replacement():
    return {"message": "Replacement"}


greet = replacement


print(route.endpoint is registered)
print(route.endpoint is greet)

print(route.endpoint())
print(greet())
\end{dhkcode}

예상 출력:

\begin{dhkzsh}
True
False
{'message': 'Original'}
{'message': 'Replacement'}
\end{dhkzsh}

왜 이런 결과가 발생할까?

처음 \zsh{@app.get("/hello")}가 적용될 때 FastAPI는 원본 \zsh{greet} 함수 객체를 등록한다.

그 이후 다음 문장이 실행된다.

\begin{dhkcode}
greet = replacement
\end{dhkcode}

이 문장은 모듈 이름 \zsh{greet}을 다른 함수 객체에 재바인딩한다.

하지만 이미 등록된 \zsh{route.endpoint}의 참조는 변경되지 않ㄴ흔다.

따라서 \zsh{FastAPI}가 등록한 엔드포인트와 모듈 이름 \zsh{greet}이 서로 다른 객체를 참조하게 된다.

이것은 \zsh{A.1}에서 학습한 이름 바인딩과 객체 동일성의 원리가 그대로 적용된 결과다.

여기서 한 가지 예외적인 상황도 고려해야 한다.

등록된 원본 함수의 본문에서 모듈 전역 이름을 다시 조회한다면, 그 전역 이름의 재바인딩이 함수 실행에 영향을 줄 수 있다.

즉, 등록된 함수 객체의 동일성이 유지된다는 사실이 함수가 사용하는 모든 외부 상태까지 고정된다는 뜻은 아니다.

\subsubsection{등록형 데코레이터의 경계 조건}

등록형 데코레이터는 함수 정의 시점에 외부 상태를 변경하므로 몇 가지 경계 조건이 존재한다.

\begin{dhktable}
|상황|	발생할 수 있는 결과|	판단 기준|
|:---|:---|:---|
|엔드포인트가 정의된 모듈이 실행되지 않음|	등록 데코레이터도 실행되지 않음	|모듈 실행 여부|
|같은 경로와 HTTP 메서드를 중복 등록|	중복 라우트가 생성될 수 있음|	경로 등록 상태와 매칭 순서|
|등록 후 함수 이름 재바인딩|	등록된 엔드포인트 참조는 그대로 유지	|객체 동일성|
|등록 전에 wrapper 적용	|wrapper가 엔드포인트로 등록됨	|데코레이터 적용 순서|
|등록 후 wrapper 적용	|원본 함수가 등록될 수 있음|	데코레이터 적용 순서|
|@wraps 없이 wrapper 등록|	시그니처 분석이 잘못될 수 있음|	__wrapped__와 실제 시그니처|
|비동기 함수를 동기 wrapper로 감쌈	|코루틴 객체가 정상적으로 처리되지 않을 수 있음|	실제 wrapper의 실행 방식|
\end{dhktable}

특히 Python 모듈을 import하면 해당 모듈의 최상위 코드가 실행될 수 있다는 사실을 기억해야 한다.

엔드포인트 정의문도 모듈 최상위 코드에 있다면 그 과정에서 실행된다.

따라서 FastAPI의 경로 등록은 일반적인 Python 모듈 실행 및 함수 정의 규칙 위에서 이루어지는 작업이다.

\subsubsection{지금까지의 Python 개념과 연결하기}

이번 학습은 새로운 문법을 학습했다기보다 이미 학습한 객체 모델을 실제 프레임워크에 적용하였다.

\begin{dhktable}
|이전 개념|	FastAPI에서의 연결|
|:---|:---|
|함수 객체와 이름 바인딩 — A.1	|엔드포인트 함수 객체 생성과 최종 이름 바인딩|
|렉시컬 스코프 — A.2~A.3	|등록 데코레이터가 경로 설정을 참조|
|클로저 — A.4~A.5|	app.get(...)이 반환하는 데코레이터의 설정 유지|
|데코레이터 변환식 — A.6	|func = decorator(func)|
|Wrapper와 반환 — A.7	|HTTP 처리에 wrapper를 적용하는 경우|
|functools.wraps — A.8|	엔드포인트 시그니처 추적|
|데코레이터 팩토리 — A.9	|@app.get("/path")|
|다중 데코레이터 순서 — A.10|	등록되는 함수 객체의 결정|
|메서드 바인딩 — A.11	|app.get이 FastAPI 인스턴스에 바인딩된 메서드|
|함수 호출과 인수 전달|	요청 데이터로부터 엔드포인트 호출 인수 구성|
\end{dhktable}

여기서 FastAPI가 Python의 기본 규칙에 추가한 기능도 구별해야 한다.

Python 자체는 함수 객체를 생성하고 호출하는 기능을 제공한다.

FastAPI는 그 함수 객체를 HTTP 경로에 등록하고, 요청을 분석하고, 의존성을 해결하고, 응답을 생성하는 기능을 제공한다.

즉, FastAPI는 Python의 함수 객체 모델을 변경하는 것이 아니라, 그 모델을 이용하여 HTTP 요청 처리 시스템을 구성하는 프레임워크다.

\subsubsection{핵심 정리와 판단 기준}

이번 학습의 전체 실행 구조를 정리해 보자.

\begin{dhksmry}[centered]
Python 모듈 실행
\da
FastAPI 인스턴스 생성
\da   
app.get("/path") 평가
\da   
등록 데코레이터 함수 객체 반환
\da   
엔드포인트 함수 객체 생성
\da   
데코레이터 적용
\da   
FastAPI 라우터에 함수 객체 등록
\da   
등록 데코레이터의 반환값을 이름에 바인딩
\da
이후 HTTP 요청 발생
\da   
요청 경로 및 HTTP 메서드 매칭
\da   
매개변수 검증·변환·의존성 해결
\da   
등록된 엔드포인트 함수 호출
\da   
반환값을 HTTP 응답으로 구성
\end{dhksmry}

이 구조를 분석할 때는 다음 다섯 가지를 확인하면 된다.

\begin{dhknote}
\begin{dhksmry}[centered]
첫째, \zsh{@app.get(...)}은 언제 평가되는가? 
\end{dhksmry}
\end{dhknote}

함수 정의문이 실행되는 과정에서 평가된다. 

HTTP 요청이 들어올 때마다 평가되는 것이 아니다.

\begin{dhknote}
\begin{dhksmry}[centered]
둘째, 무엇이 FastAPI에 등록되는가?
\end{dhksmry}
\end{dhknote}
데코레이터가 적용되는 시점에 전달받은 함수 객체가 등록된다. 

다른 데코레이터의 적용 순서에 따라 원본 함수일 수도 있고 wrapper일 수도 있다.

\begin{dhknote}
\begin{dhksmry}[centered]
셋째, 등록 이후 함수 이름과 엔드포인트 참조는 같은 객체인가?
\end{dhksmry}
\end{dhknote}

기본적인 \zsh{@app.get}만 사용하면 동일하지만, 추가 데코레이터나 이름 재바인딩이 있으면 달라질 수 있다.

\begin{dhknote}
\begin{dhksmry}[centered]
넷째, HTTP 요청에서 누가 함수를 호출하는가? 
\end{dhksmry}
\end{dhknote}
FastAPI의 요청 처리 계층이 매개변수를 준비하고 엔드포인트를 호출한다. 

따라서 직접 함수 호출과 동작이 다를 수 있다.

\begin{dhknote}
\begin{dhksmry}[centered]
다섯째, wrapper의 실제 시그니처와 실행 방식은 무엇인가?
\end{dhksmry}
\end{dhknote}

\zsh{@wraps}가 메타데이터를 제공하더라도 인수 전달과 동기·비동기 호출 방식까지 자동으로 보장하지는 않는다.

이번 단원의 가장 중요한 결론은 다음과 같다.

등록형 데코레이터는 함수 객체를 다른 객체로 변환하는 것보다, 함수 객체와 외부 실행 시스템 사이에 참조 관계를 형성하는 데 목적이 있다.

FastAPI의 \zsh{@app.get(...)}은 이 원리를 HTTP 라우팅에 적용한 대표적인 사례다.